\section{Payment rails and cost-data sources}\label{app:cost_data}\label{sec:landscape}

This appendix documents the per-rail cost data used in the corridor comparison of Section~\ref{sec:empirical_comparison}. Stablecoins compete for retail-payment volume in a space populated by infrastructure innovation along three dimensions: domestic versus cross-border, public versus private, and incumbent versus upgraded. The comparison uses a hand-assembled dataset of per-system fees, latencies, and regulatory features, deposited as a supplementary data file; public domestic instant-payment systems are characterised by their published fee schedules and regulatory caps, and the incumbent and fintech cross-border segments by the World Bank Remittance Prices Worldwide (RPW) panel \citep{rpw_q3_2025,beck_martinez_peria_2011}.

\paragraph{Upgraded public rails.} An upgraded public payment rail is operated under a public mandate, settles in the unit of account at par, and clears at zero or near-zero marginal cost to the end user within seconds. Domestically such rails are now widespread and free or near-free at the point of use: Brazil's PIX, India's UPI, the U.S. FedNow Service, the eurozone's SCT Inst (capped at the standard credit-transfer level by Regulation (EU) 2024/886), the United Kingdom's Faster Payments, and Mexico's SPEI \citep{frost_etal_2024_fast_payments,duarte_etal_2022_pix}. The operative fact for the comparison is that an upgraded public domestic rail clears at zero marginal cost to the sender. The cross-border public equivalent does not yet exist at scale but is being built, principally through BIS Project Nexus; its projected effect is taken up in Section~\ref{sec:discussion}.

\paragraph{Incumbent and fintech rails.} The dominant retail cross-border infrastructure is still SWIFT correspondent banking, whose friction stack \citep{cpmi_2018_cross_border_retail} leaves end-user costs high: a sending-bank fee, one or more intermediary fees, and an FX markup of one to four percent, summing to roughly 700 basis points on a USD 400 remittance and about 500 at USD 1{,}000. The non-bank segment is occupied by fintechs (Wise, Remitly, WorldRemit) and money-transfer operators, materially cheaper on the major corridors. The per-corridor fees used in the comparison are the RPW Q3 2025 means at the USD 500 tier; the panel reports a G20 average of 3.66 percent for sending USD 500, against the 2027 targets of 3 percent on the most expensive corridor and 1 percent globally \citep{rpw_q3_2025}.

\paragraph{Stablecoin rails.} Stablecoins are the most recent entrant, with real-economy use growing fastest in Latin America and Sub-Saharan Africa \citep{chainalysis_2025}; USDT on Tron carries the majority of retail stablecoin volume. The USDT-on-Tron fee is the most consequential rail cost for the comparison: after Tron network proposal \#104 (2025-08-29) it is on the order of USD 1.50 per transfer, about 37 basis points at USD 400 and 375 at USD 40.\footnote{At the TRX price prevailing in Q1 2026, approximately USD 1.92 for a first transfer to a new wallet, USD 1.42 for a repeat transfer, and about USD 0.50 through an energy-rental service such as TronSave \citep{tronsave_blog_2026}. USDT on Ethereum runs USD 0.50 to 7 and spikes above USD 30 in congestion; USDC on Solana costs a fraction of a cent.} These on-chain fees are observable; the two stablecoin-specific risks they do not capture are priced by the global game of Section~\ref{sec:risk_factors}. The on-ramp and off-ramp platform fees enter the use-case totals of Section~\ref{sec:empirical_comparison} as a dedicated column, because the payment circuit a retail user actually runs begins and ends in fiat: the sender converts fiat into the stablecoin at an on-ramp, and the recipient converts it back into local fiat at an off-ramp, so both legs are part of the cost of using the instrument and neither nets out against the bank or fintech alternatives, whose quoted fees are already all-in. The fees are read from posted retail schedules: spot taker fees at the major exchanges run from about 10 to 60 basis points at retail volume tiers, instant-conversion and card-funded products charge materially more, and fixed withdrawal charges add a size-dependent component. We take 25 basis points per leg as the central retail value, within the 10-to-50-basis-point range, and report the 50-basis-point round-trip in the use-case tables; the per-leg range is carried as the sensitivity reported in Section~\ref{sec:empirical_comparison}. The on-ramp leg in a developed-market funding currency (US dollar or euro) is typically at the lower end of the range and the emerging-market off-ramp leg at the upper end, so the symmetric 25-and-25 split is a conservative midpoint for the round-trip rather than a claim about either leg in isolation.
